\section{Related Work}

\paragraph{Cloth Manipulation.}
A survey on robotic cloth manipulation~\cite{longhini2024unfolding} establishes that depth and point cloud representations generalize significantly better than RGB across diverse garment appearances, because they encode geometry rather than color or texture. Methods relying on RGB inputs consistently overfit to the training distribution, while 3D representations transfer more robustly to unseen garments.

SpeedFolding~\cite{avigal2022speedfolding} uses depth and grayscale rather than RGB, further validating geometry-based representations for garment tasks. However, its folding step relies on hardcoded heuristics rather than a learned policy, and it fails on unseen garment shapes because its smoothing classifier is shape-specific and cannot generalize without retraining. FabricFlowNet~\cite{weng2022fabricflownet} leverages optical flow between current and goal-state images to predict bimanual manipulation actions, enabling goal-conditioned folding that generalizes across diverse cloth shapes. It represents an alternative visual representation direction to ours, using dense 2D flow rather than 3D point clouds.

\paragraph{Garment Folding with Learned Policies.}
UniFolding~\cite{unifolding2023} is the most directly related prior work. It uses masked partial point clouds as policy input combined with NOCS for action prediction, achieving strong cross-garment generalization on bimanual folding tasks. Their key insight is that predicting actions in a canonical coordinate space rather than raw physical space allows geometric features learned on training garments to transfer to unseen ones regardless of size or style differences. However, UniFolding uses a pick-and-place primitive action space requiring approximately 5,000 demonstrations, both of which are incompatible with our setup that requires continuous joint trajectory output and provides only 250 demonstrations per garment category. BiFold~\cite{bifold2025} enhances vision-language-action models by integrating temporal context through history buffers, allowing robots to disambiguate cloth states and recover from failures caused by self-occlusions or crumpled configurations.

\paragraph{Imitation Learning Policies.}
ACT~\cite{zhao2023act} uses RGB images with a ResNet18 encoder and has no explicit generalization mechanism. Its authors explicitly acknowledge that better perception is needed for deformable objects after their policy failed on tasks involving cloth and flexible bags. Diffusion Policy~\cite{chi2023diffusionpolicy} also uses RGB with ResNet18, but its encoder is explicitly modular and validated for swapping in ablation studies. It has been directly tested on shirt folding with approximately 284 demonstrations, closely matching our setup, and handles multimodal demonstrations naturally through stochastic diffusion. SmolVLA~\cite{smolvla2025} uses RGB with a SigLIP backbone pretrained on OCR tasks. Its vision encoder is deeply coupled to the language model backbone and the authors explicitly acknowledge long-horizon tasks as a weakness, which is directly relevant to garment folding.

\paragraph{3D Visual Representations for Policy Learning.}
DP3~\cite{ze2024dp3} integrates point cloud representations directly into a diffusion policy framework using a lightweight MLP encoder. Across 72 simulation tasks it outperforms standard Diffusion Policy by 24.2\% on average. In generalization experiments it achieves strong appearance and instance generalization while Diffusion Policy fails completely, demonstrating that the point cloud representation is the source of the generalization advantage. Critically, DP3's ablations show that a simple three-layer MLP encoder outperforms complex architectures including PointNet++~\cite{qi2017pointnet2} (78.3 vs.\ 2.2 average success rate), meaning the key is the 3D modality itself rather than encoder complexity. However, DP3's generalization advantage has only been demonstrated on rigid object tasks. Whether it transfers to deformable garment folding across unseen styles is the open question our project addresses.

\paragraph{Normalized Object Coordinate Space.}
NOCS~\cite{wang2019normalized} was originally proposed for category-level 6D object pose estimation, mapping object instances into a shared canonical space to enable generalization across instances of the same category. Our work applies this principle to manipulation policy learning, using garment keypoints as canonical anchors to normalize the geometric representation before policy conditioning.
